%% kapa: exam=1 round=28 year=2017 period=1 form=A subject=real_estate subject_ko=부동산학원론 seq=39 number=119
%% source: Q-Net 감정평가사 기출문제 (cst00309) artlSeq=5199603 file=extracted/1008/5199603/2017년 제28회 감정평가사 1차시험 문제지/1교시 A형.pdf
%% conversion: kordoc 4.24.51 -> tools/gen1.py (auto); review=proofread-sonnet; flags=none
\begin{question}{119}
부동산 경기변동과 중개활동에 관한 설명으로 옳지 않은 것은?
\begin{choices}
  \item 하향시장의 경우 종전의 거래사례 가격은 새로운 매매활동에 있어 가격 설정의 상한선이 되는 경향이 있다.
  \item 상향시장에서 매도자는 가격상승을 기대하여 거래의 성립을 미루려는 반면, 매수자는 거래성립을 앞당기려 하는 경향이 있다.
  \item 중개물건의뢰의 접수와 관련하여 안정기의 경우 공인중개사는 매각의뢰와 매입의뢰의 수집이 다 같이 중요하다.
  \item 실수요 증가에 의한 공급부족이 발생하는 경우 공인중개사는 매수자를 확보해 두려는 경향을 보인다.
  \item 일반적으로 부동산경기는 일반경기에 비하여 경기의 변동폭이 큰 경향이 있다.
\end{choices}
\end{question}
\answer{4}
